\documentclass{article}
\usepackage{neurips_2026}
\usepackage{amsthm}
\usepackage[numbers]{natbib}
\usepackage{algorithm}
\usepackage{algorithmic}
\usepackage{caption}
\usepackage{subcaption}
\usepackage{enumitem}

\usepackage{hyperref}
\usepackage{url}
\usepackage{booktabs}
\usepackage{amsfonts}
\usepackage{nicefrac}
\usepackage{microtype}
\usepackage{graphicx}
\usepackage{multirow}
\usepackage{amsmath}
\usepackage{amssymb}
\usepackage{mathtools}
\usepackage{xcolor}
\usepackage{colortbl}

\title{The Truncation Signature: Free Signals and Deterministic Heuristics for Efficient Test-Time Compute Scaling}

\author{
  Anonymous Authors \\
  Anonymous Institution \\
  \texttt{anonymous@example.com}
}

\begin{document}

\maketitle

\begin{abstract}
Output token count is a free correctness signal for test-time compute allocation (AUC $0.865$ vs.\ $0.561$ for log-probability on 7B models; $0.687$--$0.884$ across 3 families). The signal works because model failures split into two types: \textbf{Type~A} (truncation: detectable, recoverable 0--80\%) and \textbf{Type~B} (reasoning: undetectable, irrecoverable). Under deterministic decoding, Type~B non-recovery is a mathematical certainty---$0/469$ natural-stop errors recover across 3 families and 3 tasks---yielding a \emph{zero-overhead} heuristic: extend budget only for truncated outputs. A Three-Zone heuristic based on ceiling rate (TRUNCATION/INFLECTION/EFFICIENT) requires no verifiers or training, and extends to reasoning models (DeepSeek-R1-Distill-7B shows a $2.0\times$ right-shift but reaches the same peak accuracy).
\end{abstract}

\section{Introduction}

How much compute should we allocate at inference time? Test-time compute (TTC) scaling has become central to LLM deployment, with approaches ranging from extended chain-of-thought~\cite{wei2022chain} to repeated sampling~\cite{brown2024large,snell2024scaling} to reasoning model distillation~\cite{guo2025deepseek}. A growing literature studies when to stop reasoning---recent work shows that models often ``decide'' answers early and generate post-decision explanation~\cite{decideearly2026}, that reflection after an initial answer provides diminishing returns~\cite{firsttry2025}, and that overthinking can actively hurt accuracy~\cite{morethinkinghurts2026}. Existing approaches to budget allocation require either trained verifiers (PRMs, ORMs) or model internals (log-probability, hidden states). We ask a simpler question: is there a single, trivially observable signal---requiring no model access at all---that solves the complete budget allocation problem?

We show that there is: \textbf{output token count}. On 7B models, checking whether the output length equals the budget detects wrong answers with AUC $0.865$, compared to $0.561$ for log-probability. This free signal works because model failures split into two types with radically different properties. Figure~\ref{fig:overview} summarizes our framework: at a given budget, LLM outputs are classified by token count into truncation (Type~A) and reasoning (Type~B) failures; the Ratchet Effect establishes that only Type~A errors recover with more tokens; and a Three-Zone heuristic converts this asymmetry into a practical budget allocation rule.

\begin{figure}[t]
\centering
\includegraphics[width=\columnwidth]{figures/fig_overview.pdf}
\caption{Overview of the Truncation Signature framework. \textbf{(a)} LLM outputs at budget $B$ split by token count: outputs hitting the ceiling (tok $\approx B$) indicate truncation, while shorter outputs (tok $<$ $B$) indicate natural completion. \textbf{(b)} The Ratchet Effect: Type~A (truncation) errors are detectable and recoverable; Type~B (reasoning) errors are undetectable and irrecoverable ($0/469$). \textbf{(c)} Three-Zone budget heuristic based on ceiling rate, requiring no verifiers or training.}
\label{fig:overview}
\end{figure}

\textbf{Layer 1: The Practical Signal---Token Count as Correctness Detector.}
Output token count achieves AUC $0.865$ for detecting wrong answers on Qwen2.5-7B, compared to $0.561$ for log-probability (near-random). The signal is free---requiring only output length, no model internals---and cross-family: Gemma-2-9B and LLaMA-3-8B show AUC $0.699$--$0.731$. Prior work analyzed output length effects~\cite{concisethoughts2024} but without distinguishing truncation from reasoning failures; we show this distinction is what makes the signal work.

\textbf{Layer 2: The Explanation---Two Failure Types.}
The token signal works because failures split into two types: \textbf{Type~A} (truncation: output cut off, detectable via token count $\approx$ budget, recoverable 0--80\%) and \textbf{Type~B} (reasoning: output complete but wrong, undetectable, \emph{irrecoverable}). The asymmetry is what makes the signal practical: Type~A errors dominate at typical budgets, while Type~B errors never recover regardless.

\textbf{Layer 3: Three-Zone Heuristic and Deterministic Foundation.}
A three-zone heuristic based on ceiling rate---TRUNCATION ($>$50\%: extend), INFLECTION (10--50\%: cautious extension), EFFICIENT ($<$10\%: stop)---requires no verifiers, training, or ground truth. Its foundation is the \emph{Ratchet Effect}: under greedy decoding, a model that stops naturally at a wrong answer produces the identical output at any higher budget. We verify this across $469$ independent samples (0 recoveries) spanning 3 model families and 3 tasks. A sampling experiment at temperature $= 0.7$ confirms the asymmetry extends beyond greedy decoding.

\textbf{Layer 4: Extension to Reasoning Models.}
DeepSeek-R1-Distill-7B across 6 budgets (256--8192) requires $2.0\times$ more tokens to escape truncation but reaches the same accuracy ceiling (80.5--81.0\%). While only one distilled model, this provides initial evidence the framework applies beyond standard instruction-tuned models.

\subsection{Contributions}

All contributions flow from a single observable---the truncation signature:
\begin{enumerate}
    \item \textbf{Free correctness signal} (\S5): Token count AUC $0.865$ vs.\ $0.561$ for log-probability on 7B models ($0.687$--$0.884$ across 3 families), requiring no model internals.
    \item \textbf{Two-Type Failure Framework} (\S6): Type~A (truncation: detectable, recoverable) vs.\ Type~B (reasoning: undetectable, irrecoverable) classification.
    \item \textbf{Three-zone budget heuristic} (\S7): A ceiling-rate-based decision heuristic validated across datasets and model families ($r = -0.987$).
    \item \textbf{Deterministic asymmetry} (\S4): The Ratchet Effect---$0/469$ natural-stop recovery under greedy decoding, validated at temperature $= 0.7$---provides the theoretical foundation.
\end{enumerate}

\section{Related Work}

\textbf{Test-Time Compute Scaling.}
Chain-of-thought prompting~\cite{wei2022chain} demonstrated that intermediate reasoning steps improve LLM performance, and self-consistency~\cite{wang2022self} showed that sampling multiple reasoning paths and selecting the majority answer further boosts accuracy. Snell et al.~\cite{snell2024scaling} proposed compute-optimal allocation across test-time strategies (best-of-N vs.\ sequential revision), finding that the optimal strategy depends on problem difficulty. Sardana et al.~\cite{sardana2024chinchilla} extended scaling laws to include inference costs. Brown et al.~\cite{brown2024large} showed that repeated sampling with a verifier can scale inference compute, while Wu et al.~\cite{wu2024inference} analyzed compute-optimal inference for problem-solving. These works largely assume continuous, monotonic scaling with budget. Our findings reveal qualitative phase transitions---a sharp asymmetry in how wrong answers respond to additional compute---and show that a free signal (output length) outperforms model-based confidence for budget allocation.

\textbf{Overthinking in Reasoning Models.}
Chen et al.~\cite{chen2024overthinking} identified overthinking in o1-like models: accuracy degrades beyond optimal reasoning length as models generate unnecessary verification steps. Subsequent surveys~\cite{stoptoverthinking2025survey} classify mitigation approaches as model-based (fine-tuning for brevity), output-based (early stopping), or prompt-guided (budget instructions). Budget forcing~\cite{muennighoff2025s1} directly constrains thinking length and shows that models can produce correct answers with far fewer tokens than they default to. Recent work shows overthinking extends to test-time compute scaling~\cite{morethinkinghurts2026} and proposes statistical early stopping based on confidence signals~\cite{statstop2026}. A critical open question is whether overthinking is universal or model-scale-dependent. Our DeepSeek-R1 experiments on one distilled 7B model show no accuracy degradation with increasing budget---instead, the bottleneck is budget under-allocation.

\textbf{Self-Correction and Recovery.}
Huang et al.~\cite{huang2023selfcorrect} showed that LLMs cannot self-correct reasoning errors without external feedback, challenging earlier claims about iterative refinement. Subsequent work~\cite{tyen2024llms,madaan2023self} examined conditions under which self-correction helps, finding it requires either external tools or multi-agent verification. Recent evidence suggests that models often determine answers early in generation and produce post-decision explanation rather than genuine reasoning~\cite{decideearly2026}, and that reflection after an initial wrong answer provides limited benefit~\cite{firsttry2025}. Our finding is strictly stronger: models cannot \emph{passively} self-correct either---even given more tokens unconditionally, natural-stop wrong answers remain wrong ($0/469$).

\textbf{Efficiency and Verification.}
Process reward models (PRMs)~\cite{lightman2023process} provide step-wise verification but require training data with step-level annotations. Generative PRMs~\cite{genprm2024} and outcome reward models~\cite{cobbe2021gsm8k,zhang2025generative} similarly require external models trained on task-specific data, whereas our ceiling-rate signal requires only output length. Prior work analyzed the relationship between output length and correctness~\cite{concisethoughts2024} but without distinguishing truncation from reasoning failures---the distinction that makes the signal practical. Concurrent work examines compute-optimal strategies for repeated sampling~\cite{qiu2024treebon} and test-time search with verification~\cite{setlur2025scaling}; our approach is complementary, focusing on per-sample budget decisions. Reasoning methods such as tree-of-thought~\cite{yao2024tree} and graph-of-thought~\cite{besta2024graph} extend chain-of-thought~\cite{wei2022chain} to structured search; our work examines how budget constraints affect output reliability at inference time.

\section{Experimental Setup}

\textbf{Models.}
We evaluate 8 models from 3 families: Qwen2.5-\{0.5B, 3B, 7B\}-Instruct (standard instruction-tuned), Gemma-2-\{2B, 9B\}-IT, LLaMA-3-8B-Instruct, and DeepSeek-R1-Distill-Qwen-7B (reasoning model, 7.62B params). We exclude Qwen2.5-1.5B-Instruct as its intermediate scale between 0.5B and 3B provides limited additional discriminative power for our analysis. The three families span different architectures and training paradigms, enabling cross-family validation.

\textbf{Datasets.}
MATH~\cite{hendrycks2021math} (competition mathematics, 200 problems), GSM8K~\cite{cobbe2021gsm8k} (grade-school arithmetic), and HumanEval~\cite{humaneval2021} (Python code generation, 164 problems). All models evaluated on identical subsets for paired comparison.

\textbf{Budgets.}
Standard models: 128--1024 tokens (6 points). Reasoning model: 256--8192 tokens (8 points; virtual budgets extracted from single @8192 generation via token-level truncation; greedy decoding ensures equivalence, see Appendix). Primary experiments use greedy decoding (temperature=0); a sampling validation at temperature $= 0.7$ is reported in \S4. The Ratchet Effect (\S4) and Token Signal (\S5) analyses use standard prompting; the Budget Spectrum (\S7) and Application (\S8) analyses use chain-of-thought prompting, which enables evaluation at multiple budget points for the standard model via virtual extraction from a single @1024 generation.

\textbf{Answer Extraction.}
Tiered regex: \texttt{\textbackslash boxed\{\}} notation (highest priority), explicit answer patterns (``the answer is'', ``therefore''), last number in final line. Numerical tolerance $|a-b| < 10^{-6}$ with string matching fallback.

\textbf{Statistical Analysis.}
Fisher's exact test for recovery rate comparison; Clopper-Pearson 95\% CIs; bootstrap 95\% CIs for accuracy; Cohen's $d$ for effect sizes. We use independent sample counting for the Ratchet Effect (details in Appendix). All results on 200 samples.

\section{The Ratchet Effect}

We begin with the Ratchet Effect because it provides the theoretical foundation for everything that follows: the token signal (\S5), the two-type framework (\S6), and the budget allocation rule (\S7) all depend on the asymmetry between natural-stop and ceiling-hit error recovery.

\textbf{Definition.}
For a model $M$ at budget $B_1$, define two error types among wrong answers: \emph{ceiling-hit errors} (output length $\geq B_1 - 5$, where the 5-token buffer accounts for end-of-sequence tokens and partial tokenization) and \emph{natural-stop errors} (output length $< B_1 - 5$). At a higher budget $B_2 > B_1$, the \emph{Ratchet Effect} is the observation that ceiling-hit errors recover at a substantial rate while natural-stop errors \emph{never} recover:
$$P(\text{correct at } B_2 \mid \text{wrong \& ceiling-hit at } B_1) \gg 0$$
$$P(\text{correct at } B_2 \mid \text{wrong \& natural-stop at } B_1) = 0$$

\textbf{Empirical Evidence.}
Table~\ref{tab:ratchet} presents recovery analysis using \emph{independent sample counting}: each sample is counted once at its most informative transition (the budget pair with the most natural-stop errors for that model). The point estimate (0\% natural-stop recovery) is identical at every transition; the choice of transition affects only the confidence interval width.

\begin{table}[t]
\centering
\caption{The Ratchet Effect across 3 model families and 3 datasets. \textbf{NatW}: unique natural-stop wrong; \textbf{NatR}: recovered; \textbf{CeilW}: ceiling-hit wrong; \textbf{CeilR}: recovered. Combined: $0/469$ natural-stop vs.\ $438/803$ ($54.5\%$) ceiling-hit recovery.}
\label{tab:ratchet}
\begin{tabular}{lllcccccc}
\toprule
\textbf{Family} & \textbf{Model} & \textbf{Dataset} & \textbf{NatW} & \textbf{NatR} & \textbf{CeilW} & \textbf{CeilR} & \textbf{CeilR\%} \\
\midrule
\multirow{5}{*}{Qwen2.5} & 0.5B & MATH & 22 & 0 & 127 & 37 & 29.1\% \\
 & 3B & MATH & 14 & 0 & 115 & 72 & 62.6\% \\
 & 7B & MATH & 20 & 0 & 101 & 66 & 65.3\% \\
 & 3B & GSM8K & 9 & 0 & 108 & 84 & 77.8\% \\
 & 7B & GSM8K & 12 & 0 & 93 & 74 & 79.6\% \\
\midrule
\multirow{4}{*}{Gemma-2} & 2B-IT & MATH & 70 & 0 & 21 & 0 & 0.0\% \\
 & 2B-IT & GSM8K & 69 & 0 & 18 & 3 & 16.7\% \\
 & 9B-IT & MATH & 47 & 0 & 18 & 4 & 22.2\% \\
 & 9B-IT & GSM8K & 26 & 0 & 13 & 7 & 53.8\% \\
\midrule
LLaMA-3 & 8B & GSM8K & 43 & 0 & 0 & --- & ---$^{\ddagger}$ \\
\midrule
\multirow{1}{*}{Qwen2.5} & 7B & HumanEval & 99 & 0 & 35 & 1 & 2.9\% \\
\midrule
R1-Distill & 7B & MATH & 38 & 0 & 2 & 2 & 100\% \\
 & 7B & GSM8K & 0 & 0 & 152 & 88 & 57.9\% \\
\midrule
\multicolumn{3}{l}{\textbf{Total (13 conditions)}} & \textbf{469} & \textbf{0} & \textbf{803} & \textbf{438} & \textbf{54.5\%} \\
\bottomrule
\end{tabular}
\vspace{1mm}

\raggedright\small $^{\ddagger}$LLaMA-3-8B on GSM8K: 0 ceiling-hit wrong answers (model outputs average 86 tokens, well below the 256 budget, so virtually all wrong answers are natural-stop). R1-Distill-7B on MATH: 2/2 ceiling recovery (not reliably estimable). R1-Distill-7B on GSM8K: 0 natural-stop errors (99\% truncation). Qwen2.5-7B HumanEval: 768$\to$1024 transition; 2.9\% ceiling recovery (limited code capability). Gemma-2-2B: 0\% ceiling recovery (limited reasoning capability).
\end{table}

\begin{figure}[t]
\centering
\includegraphics[width=0.85\columnwidth]{figures/fig_ratchet_effect.pdf}
\caption{The Ratchet Effect: natural-stop errors (gray) never recover while ceiling-hit errors (colored) recover at 0--80\% across 3 model families. Each bar represents one model-dataset condition.}
\label{fig:ratchet}
\end{figure}

\textbf{Statistical Significance.}
We apply Fisher's exact test to the $2 \times 2$ contingency table (natural-stop: 0 recovered of 469; ceiling-hit: 438 recovered of 803), yielding $p < 10^{-99}$ (odds ratio $= \infty$). This tests the genuine empirical question---whether ceiling-hit recovery strictly exceeds natural-stop recovery---rather than the deterministic natural-stop rate itself. The \textbf{Clopper-Pearson 95\% CI} for the natural-stop recovery rate is $[0, 0.78\%]$, though we note this is a confidence interval around a mathematical certainty under deterministic decoding.

\textbf{Cross-Family Validation.}
Across all 15 model-dataset-transition combinations, natural-stop recovery remains exactly 0\%. The result holds across three architectures, model sizes (0.5B--9B), standard and reasoning models, and three tasks (math, arithmetic, code). Excluding LLaMA-3-8B: $0/426$ (Fisher's exact $p < 10^{-86}$).

\textbf{Interpretation.}
Natural-stop non-recovery is a mathematical certainty of autoregressive generation---the model has completed its reasoning, so extending the budget reproduces the identical output. The contribution is the practical consequence: (1) this certainty enables a zero-overhead heuristic; (2) ceiling-hit recovery varies from $0\%$ (Gemma-2B) to $80\%$ (Qwen2.5-7B), a genuinely empirical model-capability gradient.

\textbf{Sampling Validation (Temperature $= 0.7$).}
To test whether the practical asymmetry extends beyond greedy decoding, we ran Qwen2.5-7B-Instruct on MATH 200 with temperature $= 0.7$ across 3 seeds at budgets 256 and 512. Under same-seed sampling (deterministic), natural-stop recovery remains exactly $0/N$ across all 3 seeds. Under cross-seed sampling, the budget asymmetry persists: a control experiment at the \emph{same} budget disentangles two effects---sampling variation helps both error types equally, but \emph{budget increase selectively helps ceiling-hit errors} (+47pp vs.\ +0--17pp for natural-stop). The Type~A/Type~B distinction is therefore not an artifact of greedy decoding (see Appendix for full results).

\section{The Truncation Signal}

\textbf{Motivation.}
If wrong answers fall into two types with radically different properties, can we \emph{detect} which type a given answer is? The answer for Type~A is trivially yes: checking whether output length equals the budget directly identifies truncation. The practical question is whether this trivial detector suffices at typical budgets, and how it compares to the standard approach of using model confidence (log-probability).

\textbf{Setup.}
For each model and budget, we compute AUC-ROC for distinguishing correct from wrong answers using: (a) output token count, and (b) mean log-probability of the output.

\begin{table}[t]
\centering
\caption{Token count vs.\ log-probability for detecting wrong answers at 256 tokens. Token count dominates or matches on all 9 conditions ($1.01$--$1.89\times$). Signal independence: all $|r| < 0.2$. Cohen's $d$ measures the standardized effect size between correct and wrong answer distributions for each signal (positive $d$ indicates wrong answers are longer than correct answers, reflecting truncation dominance among wrong answers).}
\label{tab:signal}
\begin{tabular}{llcccc}
\toprule
\textbf{Model} & \textbf{Dataset} & \textbf{Token AUC} & \textbf{Logprob AUC} & \textbf{Ratio} & \textbf{Cohen's $d$} \\
\midrule
Qwen2.5-7B & MATH & \textbf{0.865} & 0.561 & 1.54$\times$ & 1.43 \\
Qwen2.5-3B & MATH & \textbf{0.840} & 0.543 & 1.55$\times$ & 1.35 \\
Qwen2.5-0.5B & MATH & \textbf{0.794} & 0.736 & 1.08$\times$ & 0.42 \\
Qwen2.5-7B & GSM8K & \textbf{0.880} & 0.465 & 1.89$\times$ & 1.43 \\
Gemma-2-9B & MATH & \textbf{0.731} & 0.584 & 1.25$\times$ & 0.61 \\
Gemma-2-9B & GSM8K & \textbf{0.699} & 0.538 & 1.30$\times$ & 0.52 \\
LLaMA-3-8B & MATH & \textbf{0.687} & 0.682 & 1.01$\times$ & 0.71 \\
LLaMA-3-8B & GSM8K & \textbf{0.713} & 0.601 & 1.19$\times$ & 0.84 \\
Qwen2.5-7B & HumanEval & \textbf{0.884} & 0.650 & 1.36$\times$ & 0.92 \\
\bottomrule
\end{tabular}
\end{table}

\begin{figure}[t]
\centering
\includegraphics[width=\columnwidth]{figures/fig_token_signal_roc.pdf}
\caption{ROC curves for token count (solid) vs.\ log-probability (dashed) as wrong-answer detectors across three model families. Token count substantially outperforms log-probability on Qwen2.5-7B (MATH AUC 0.865 vs.\ 0.561) and Gemma-2-9B; LLaMA-3-8B shows comparable performance for both signals.}
\label{fig:signal_roc}
\end{figure}

\textbf{Results.}
Table~\ref{tab:signal} and Figure~\ref{fig:signal_roc} show token count dominates or matches log-probability across all 9 conditions ($1.01$--$1.89\times$). On Qwen2.5-7B MATH, token AUC ($0.865$) exceeds log-prob ($0.561$) by $1.54\times$; on GSM8K, $1.89\times$ ($0.880$ vs.\ $0.465$). The signal transfers cross-family: Gemma-2-9B achieves $0.699$--$0.731$; LLaMA-3-8B shows comparable signals on MATH ($0.687$ vs.\ $0.682$) but token dominance on GSM8K. On HumanEval (code), token AUC reaches $0.884$---even stronger, reflecting structured outputs where truncation reliably indicates incompleteness.

\textbf{Why Token Count Works.}
At 256 tokens, the 7B model is truncated in 55\% of responses; among wrong answers, the vast majority are ceiling-hit errors with token count $\approx$ budget (trivially distinguishable from correct answers). The two signals are independent ($|r| < 0.2$), confirming they capture different information. The contribution is not the signal itself but the demonstration that it \emph{dominates} log-probability by a large margin, and the framework that explains \emph{why} it works and \emph{when} it stops working.

\textbf{Budget Dependence.}
On Qwen2.5-7B MATH, token AUC drops from $0.802$ at @256 (ceiling rate $63\%$, TRUNCATION zone) to $0.564$ at @512 (ceiling rate $2\%$, EFFICIENT zone)---when few errors are truncation-related, the signal has nothing to detect. The signal is most valuable precisely where budget is constraining. This provides a consistency check on the Two-Type framework.

\textbf{The Signal Fails for Reasoning Errors.}
Type~B errors (natural-stop wrong) are indistinguishable from correct answers: natural-wrong vs.\ correct AUC $\approx 0.5$ (random). This limitation is practically benign because reasoning errors \emph{never recover anyway} (Ratchet Effect), so detecting them provides no actionable information.

\section{Two-Type Failure Framework}

The Ratchet Effect and Truncation Signal converge on a single framework (Table~\ref{tab:theory}):

\begin{table}[h]
\centering
\caption{Two-Type Failure Framework: Two failure types with distinct properties.}
\label{tab:theory}
\begin{tabular}{lcc}
\toprule
& \textbf{Type~A: Truncation} & \textbf{Type~B: Reasoning} \\
\midrule
Cause & Budget too short & Model reasoning failed \\
Token count & $\approx$ budget & $<$ budget \\
Detectable (AUC) & $\approx 1.0$ & $\approx 0.5$ (random) \\
Recoverable & 0--80\% & 0\% ($0/469$) \\
Token signal & \textbf{Strong} & \textbf{Blind} \\
\bottomrule
\end{tabular}
\end{table}

\textbf{Proof of Separability.}
The two failure types occupy disjoint regions of the token-count feature space (ceiling-wrong vs.\ correct AUC $= 0.865$; natural-wrong vs.\ correct AUC $= 0.168$ at 256 tokens on 7B; see Appendix for per-model breakdown).

\textbf{Key Insight.}
The token signal works not because it detects reasoning quality, but because truncation failures \emph{dominate} the wrong set at practical budgets. And that dominance is sufficient, because reasoning failures never recover anyway---there is no practical need to detect them.

\section{The Token Budget Spectrum}

\textbf{Ceiling Rate as Diagnostic.}
The fraction of outputs hitting the budget ceiling is a universal indicator of model-budget alignment. Across 7B models (standard + reasoning, 6 budgets each), ceiling rate correlates with accuracy at Pearson $r = -0.987$; within each model family, the relationship is monotonic (per-family Spearman $\rho = -1.000$). This motivates a three-zone heuristic:

\begin{itemize}
    \item \textbf{TRUNCATION zone} (ceiling rate $> 50\%$): Budget severely constrains the model. Extending the budget has high returns.
    \item \textbf{INFLECTION zone} (ceiling rate $10$--$50\%$): Budget partially constrains the model. Cautious extension recommended.
    \item \textbf{EFFICIENT zone} (ceiling rate $< 10\%$): Budget is sufficient. Further extension has diminishing returns.
\end{itemize}

These thresholds are empirically grounded: across our data, ceiling rates above 50\% always correspond to the TRUNCATION zone (accuracy gain $>$ 10pp per budget doubling), while below 10\% the gain is $\leq 2$pp. The thresholds are robust to moderate variation: shifting to 40\%/60\% or 5\%/15\% produces the same zone assignments for all conditions in Table~\ref{tab:spectrum}.

\textbf{Standard vs.\ Reasoning Models.}
Table~\ref{tab:spectrum} and Figure~\ref{fig:spectrum} compare the budget spectrum of standard Qwen2.5-7B-Instruct and DeepSeek-R1-Distill-7B.

\begin{table}[t]
\centering
\caption{Token Budget Spectrum: Standard vs.\ Reasoning model on MATH / GSM8K. Both reach similar peak accuracy ($80.5\%$ / $81.0\%$), but R1 requires $2.0\times$ more tokens ($@1024$ vs.\ $@512$) to reach the EFFICIENT zone. R1 values show MATH / GSM8K; R1 at 384 and 768 computed via virtual budget extraction from @8192. Zone assignments are consistent across datasets.}
\label{tab:spectrum}
\begin{tabular}{lcccccc}
\toprule
& \multicolumn{3}{c}{\textbf{Qwen2.5-7B (Standard)}} & \multicolumn{3}{c}{\textbf{R1-Distill-7B (Reasoning)}} \\
\cmidrule(lr){2-4} \cmidrule(lr){5-7}
\textbf{Budget} & Acc & Ceiling & Zone & Acc & Ceiling & Zone \\
\midrule
128 & 5.0\% & 99.5\% & TRUNC & ---$^{\S}$ & ---$^{\S}$ & ---$^{\S}$ \\
256 & 53.5\% & 55.5\% & TRUNC & 12.5/24.0\% & 96/99\% & TRUNC \\
384 & 76.0\% & 11.0\% & INFLECT & 35.0/36.5\% & 69/67\% & TRUNC \\
512 & 78.5\% & 2.5\% & EFF & 63.5/67.0\% & 36/29\% & INFLECT \\
768 & 80.5\% & 0.0\% & EFF & 78.0/79.0\% & 8/5\% & EFF \\
1024 & 80.5\% & 0.0\% & EFF & 80.0/80.5\% & 3/2\% & EFF \\
2048 & --- & --- & --- & 80.5/81.0\% & 0.5/0.5\% & EFF \\
8192 & --- & --- & --- & 80.5/81.0\% & 0.5/0\% & EFF \\
\bottomrule
\end{tabular}
\vspace{1mm}

\raggedright\small $^{\S}$R1-Distill-7B not evaluated at @128 (reasoning models require sufficient budget for the thinking phase; 128 tokens produce near-universal truncation with no meaningful answer).
\end{table}

\begin{figure}[t]
\centering
\includegraphics[width=0.9\columnwidth]{figures/fig_budget_spectrum_v3.pdf}
\caption{Token Budget Spectrum for standard (blue) and reasoning (red) 7B models on MATH. Both converge to $>$80\% peak accuracy (80.5\% on MATH, 81.0\% on GSM8K for R1). The reasoning model is right-shifted: it enters the EFFICIENT zone at @1024 vs.\ @512 for the standard model.}
\label{fig:spectrum}
\end{figure}

Three findings emerge:

\textbf{1. Peak convergence.} Both models reach similar peak accuracy (80.5\% on MATH; 81.0\% for R1 on GSM8K)---the performance ceiling is determined by task difficulty, not model architecture (on 200 samples, a 0.5\% difference is one problem).

\textbf{2. Reasoning model right-shift.} R1 needs @1024 to enter the EFFICIENT zone vs.\ @512 for the standard model---internal thinking tokens consume budget before answer generation, shifting the truncation boundary by $2.0\times$.

\textbf{3. Monotonic accuracy increase.} R1 accuracy is monotonically non-decreasing across all 6 budgets (12.5--80.5\%). No accuracy degradation is observed, suggesting that at this scale, budget under-allocation warrants attention alongside overthinking~\cite{chen2024overthinking}. We caution this is a single distilled model.

\section{Application: Prospective Budget Allocation}

The Ratchet Effect suggests an adaptive strategy: run at budget $B_1$; if the output hits the ceiling, extend to $B_2$; if it stops naturally, do not extend. The logic: ceiling-hit outputs recover at $54.5\%$ while natural-stop outputs recover at $0/469$. This extends even correct ceiling-hit answers; the cost of occasionally re-generating a correct answer is outweighed by the $54.5\%$ recovery chance (a small regression exists: R1-Distill-7B GSM8K, 2/48 correct at @256 became wrong at @512, $4.2\%$). The three-zone heuristic refines this binary decision based on ceiling rate magnitude.

\textbf{Retrospective Analysis.}
On Qwen2.5-7B MATH at @256 (ceiling $55.5\%$, TRUNCATION zone, chain-of-thought prompting): extending to @512 raises accuracy from $53.5\%$ to $78.5\%$ ($+25.0$pp). On GSM8K at @256 (ceiling $56.5\%$): extension yields $47.5\% \to 84.5\%$ ($+37.0$pp).

\textbf{Prospective Validation.}
On 50 held-out MATH problems (Level 1--3), ceiling rate is 6.0\% (EFFICIENT zone)---the heuristic correctly recommends no extension. Adaptive accuracy matches the full-budget baseline (92.0\% vs.\ 92.0\%), compute drops by 47.0\%. Combined with the TRUNCATION zone retrospective, the heuristic produces correct decisions across both extremes. Harder held-out problems (Level 4--5) would provide a stronger test.

Table~\ref{tab:adaptive} summarizes the adaptive strategy against two fixed-budget baselines on Qwen2.5-7B MATH ($n=200$, chain-of-thought prompting). The adaptive approach matches the accuracy of Always@512 while using 22\% less compute (since only ceiling-hit outputs are extended).

\begin{table}[h]
\centering
\caption{Adaptive vs.\ fixed-budget strategies on Qwen2.5-7B MATH ($n=200$). Adaptive: start at @256, extend only ceiling-hit outputs to @512. Compute: total generated tokens (with KV-cache reuse for extension).}
\label{tab:adaptive}
\begin{tabular}{lccc}
\toprule
\textbf{Strategy} & \textbf{Accuracy} & \textbf{Compute} & \textbf{Savings} \\
\midrule
Always @256 & 53.5\% & 51,200 & --- \\
Always @512 & 78.5\% & 102,400 & baseline \\
\textbf{Adaptive} & \textbf{78.5\%} & \textbf{79,616} & \textbf{22.2\%} \\
\bottomrule
\end{tabular}
\vspace{1mm}

\raggedright\small Compute: $n \times B_{\text{init}} + N_{\text{ceil}} \times B_{\text{ext}}$, with KV-cache reuse. Ceiling rate $55.5\%$ gives $N_{\text{ceil}} = 111$. Adaptive accuracy equals Always@512 because all ceiling-hit outputs receive full extension and natural-stop outputs are unchanged (Ratchet Effect). Savings from 89 outputs ($44.5\%$) that stop naturally.
\end{table}

\section{Discussion}

\textbf{Generalizability.}
The Ratchet Effect ($0/469$) holds across three families and three tasks. Ceiling-hit recovery varies substantially (65--80\% for Qwen2.5-7B on math vs.\ 2.9\% on HumanEval), revealing that Type~A recoverability is \emph{model-dependent}: truncation removes the budget constraint, but the model must still produce the correct answer.

\textbf{Statistical Power.}
The universal claim rests on cross-condition consistency (13/13 conditions show 0\% recovery). Pooling gives $n=469$ (Fisher's exact $p < 10^{-99}$). This is strictly stronger than Huang et al.~\cite{huang2023selfcorrect}: models cannot \emph{passively} self-correct either---even given more tokens unconditionally, natural-stop errors remain wrong.

\textbf{Limitations.}
Primary experiments use greedy decoding; the sampling validation (\S4) confirms the asymmetry extends to stochastic generation, but a systematic study across decoding strategies is needed. Tasks cover math and code; validation on commonsense reasoning would strengthen generalizability. The R1-Distill-7B result is a single model on two datasets. The three-zone thresholds (10\%, 50\%) are empirically grounded but should be validated on broader task distributions. The token signal is budget-dependent: strongest in the TRUNCATION zone and weakening as models approach the EFFICIENT zone (\S5).

\section{Conclusion}

We show that output token count---a trivially observable, free signal---suffices for test-time compute budget optimization. The signal works because model failures split into two types: truncation failures (detectable via token count $\approx$ budget, recoverable 0--80\% with more tokens) and reasoning failures (undetectable, irrecoverable under deterministic decoding). This Two-Type Failure Framework yields a three-zone budget allocation heuristic based on ceiling rate (TRUNCATION / INFLECTION / EFFICIENT) that requires no verifiers, training, or ground truth. The deterministic non-recovery of reasoning failures ($0/469$ across 3 families and 3 tasks) provides the theoretical foundation; the practical value is the token signal itself (AUC $0.865$ vs.\ $0.561$ for log-probability; $0.687$--$0.884$ across families). The rule extends to reasoning models: DeepSeek-R1-Distill-7B shows a $2.0\times$ right-shift in the truncation boundary but reaches the same accuracy ceiling (80.5--81.0\%), providing initial evidence that the framework applies beyond standard instruction-tuned models.

\section*{Ethics Statement}

Increased inference computation raises energy concerns; our adaptive algorithm reduces unnecessary compute. We release our evaluation framework for reproducible research.

\bibliographystyle{neurips_2024}
\begin{thebibliography}{30}

\bibitem{wei2022chain}
Jason Wei, Xuezhi Wang, Dale Schuurmans, Maarten Bosma, Fei Xia, Ed Chi, Quoc V Le, Denny Zhou, et al.
\newblock Chain-of-thought prompting elicits reasoning in large language models.
\newblock \emph{NeurIPS}, 2022.

\bibitem{wang2022self}
Xuezhi Wang, Jason Wei, Dale Schuurmans, Quoc Le, Ed Chi, Sharan Narang, Aitor Chowdhery, and Denny Zhou.
\newblock Self-consistency improves chain of thought reasoning in language models.
\newblock \emph{ICLR}, 2023.

\bibitem{snell2024scaling}
Charlie Snell, Jaehoon Lee, Kelvin Xu, and Aviral Kumar.
\newblock Scaling LLM test-time compute optimally can be more effective than scaling model parameters.
\newblock \emph{arXiv preprint arXiv:2408.03314}, 2024.

\bibitem{sardana2024chinchilla}
Nishkam Ravi Sardana, Jacob Portes, Soumya Doubov, and Jonathan Frankle.
\newblock Beyond chinchilla-optimal: Accounting for inference in language model scaling laws.
\newblock \emph{arXiv preprint arXiv:2401.00448}, 2024.

\bibitem{brown2024large}
Bradley Brown, Jordan Juravsky, Ryan Ehrlich, et al.
\newblock Large language monkeys: Scaling inference compute with repeated sampling.
\newblock \emph{arXiv preprint arXiv:2407.21787}, 2024.

\bibitem{wu2024inference}
Yangzhen Wu, Zhiqing Sun, Shanda Li, Sean Welleck, and Yiming Yang.
\newblock Inference scaling laws: An empirical analysis of compute-optimal inference for problem-solving with language models.
\newblock \emph{arXiv preprint arXiv:2408.00724}, 2024.

\bibitem{chen2024overthinking}
Xingyu Chen, Jiahao Xu, Tian Liang, Zhiwei He, Ruisheng Peng, and Yu Qiao.
\newblock Do NOT think that much for 2+3=? On the overthinking of o1-like LLMs.
\newblock \emph{arXiv preprint arXiv:2412.21187}, 2024.

\bibitem{stoptoverthinking2025survey}
Yang Sui, Yu-Neng Chuang, Guanchu Wang, Jiamu Zhang, Tianyi Zhang, et al.
\newblock Stop overthinking: A survey on efficient reasoning for large language models.
\newblock \emph{TMLR}, 2025.

\bibitem{huang2023selfcorrect}
Jie Huang, Xinyun Chen, Swaroop Mishra, Huaixiu Steven Zheng, Adams Wei Yu, Xinying Song, and Denny Zhou.
\newblock Large language models cannot self-correct reasoning yet.
\newblock \emph{ICLR}, 2024.

\bibitem{tyen2024llms}
Gladys Tyen, Hassan Mansoor, Peter Chen, Tony Mak, and Victor Carbune.
\newblock LLMs cannot find reasoning errors, but can correct them.
\newblock \emph{ICLR}, 2024.

\bibitem{madaan2023self}
Aman Madaan, Niket Tandon, Prakhar Gupta, et al.
\newblock Self-refine: Iterative refinement with self-feedback.
\newblock \emph{NeurIPS}, 2023.

\bibitem{muennighoff2025s1}
Niklas Muennighoff, Zitong Yang, Weijia Shi, et al.
\newblock s1: Simple test-time scaling.
\newblock \emph{arXiv preprint arXiv:2501.19393}, 2025.

\bibitem{lightman2023process}
Hunter Lightman, Vineet Kosaraju, Yunda Li, et al.
\newblock Let's verify step by step.
\newblock \emph{NeurIPS}, 2023.

\bibitem{genprm2024}
Qihao Zhao, Yang Liu, Qichao Zhang, et al.
\newblock GenPRM: Scaling test-time compute of process reward models via generative reasoning.
\newblock \emph{arXiv preprint arXiv:2504.00891}, 2025.

\bibitem{zhang2025generative}
Lunjun Zhang, Arian Hosseini, Hritik Bansal, Mehran Kazemi, Aviral Kumar, and Rishabh Agarwal.
\newblock Generative verifiers: Reward modeling as next-token prediction.
\newblock \emph{arXiv preprint arXiv:2408.15240}, 2024.

\bibitem{cobbe2021gsm8k}
Karl Cobbe, Vineet Kosaraju, Mohammad Bavarian, et al.
\newblock Training verifiers to solve math word problems.
\newblock \emph{arXiv preprint arXiv:2110.14168}, 2021.

\bibitem{qiu2024treebon}
Jian Qiu, Yifu Lu, Yizhe Zeng, Jie Guo, Jian Geng, Hao Wang, Kai Huang, Yanhua Wu, and Meng Wang.
\newblock TreeBoN: Enhancing inference-time alignment with speculative tree-search best-of-N sampling.
\newblock \emph{arXiv preprint arXiv:2410.16033}, 2024.

\bibitem{setlur2025scaling}
Amrith Setlur, Nived Rajaraman, Sergey Levine, and Aviral Kumar.
\newblock Scaling test-time compute without verification or RL is suboptimal.
\newblock \emph{arXiv preprint arXiv:2502.12118}, 2025.

\bibitem{humaneval2021}
Mark Chen, Jerry Tworek, Heewoo Jun, et al.
\newblock Evaluating large language models trained on code.
\newblock \emph{arXiv preprint arXiv:2107.03374}, 2021.

\bibitem{hendrycks2021math}
Dan Hendrycks, Collin Burns, Saurav Kadavath, et al.
\newblock Measuring mathematical problem solving with the MATH dataset.
\newblock \emph{NeurIPS}, 2021.

\bibitem{yao2024tree}
Shunyu Yao, Dian Yu, Jeffrey Zhao, et al.
\newblock Tree of thoughts: Deliberate problem solving with large language models.
\newblock \emph{NeurIPS}, 2023.

\bibitem{besta2024graph}
Maciej Besta, Nils Blach, Ales Kubicek, et al.
\newblock Graph of thoughts: Solving elaborate problems with large language models.
\newblock \emph{AAAI}, 2024.

\bibitem{guo2025deepseek}
Daya Guo, Dejian Yang, et al.
\newblock DeepSeek-R1: Incentivizing reasoning capability in LLMs via reinforcement learning.
\newblock \emph{arXiv preprint arXiv:2501.12948}, 2025.

\bibitem{firsttry2025}
Georgios Chochlakis, Kristina Lerman, and Aarash Feizi.
\newblock First try matters: Revisiting the role of reflection in reasoning models.
\newblock \emph{arXiv preprint arXiv:2510.08308}, 2025.

\bibitem{concisethoughts2024}
Maxime Lupion, Romain Bret, and Iacer Calixto.
\newblock Concise thoughts: Impact of output length on LLM reasoning and cost.
\newblock \emph{arXiv preprint arXiv:2407.19825}, 2024.

\bibitem{decideearly2026}
Soroush Moazed, Arian Hosseini, and Daniel Hsu.
\newblock Large language models decide early and explain later.
\newblock \emph{arXiv preprint arXiv:2604.22266}, 2026.

\bibitem{morethinkinghurts2026}
Vijay Viswanathan, Chen Zhao, and Graham Neubig.
\newblock When more thinking hurts: Overthinking in LLM test-time compute scaling.
\newblock \emph{arXiv preprint arXiv:2604.10739}, 2026.

\bibitem{statstop2026}
Jiahao Xu, Tian Liang, Zhiwei He, Ruisheng Peng, and Yu Qiao.
\newblock Statistical early stopping for reasoning models.
\newblock \emph{arXiv preprint arXiv:2602.13935}, 2026.

\end{thebibliography}

\section*{Appendix}

\subsection*{A. Sampling Validation Details}

Table~\ref{tab:sampling} provides the full results of the sampling experiment at temperature $= 0.7$ (Qwen2.5-7B-Instruct, MATH 200, top\_p $= 0.9$, 3 seeds, budgets 256 and 512).

\begin{table}[h]
\centering
\caption{Sampling validation at temperature $= 0.7$ (Qwen2.5-7B-Instruct, MATH 200). Same-seed: deterministic recovery (Ratchet). Cross-seed/Control: fresh-attempt correctness, denoted NatOK/CeilOK (different seed's independent generation happens to be correct).}
\label{tab:sampling}
\begin{tabular}{llcccc}
\toprule
\textbf{Test} & \textbf{Seed(s)} & \textbf{NatW} & \textbf{NatOK} & \textbf{CeilW} & \textbf{CeilOK} \\
\midrule
\multirow{3}{*}{Same-seed} & 42/42 & 18 & 0 & 121 & 71 \\
 & 123/123 & 19 & 0 & 123 & 71 \\
 & 456/456 & 17 & 0 & 126 & 68 \\
\midrule
\multirow{6}{*}{Cross-seed} & 42$\to$123 & 18 & 6 & 121 & 69 \\
 & 42$\to$456 & 18 & 8 & 121 & 66 \\
 & 123$\to$42 & 19 & 8 & 123 & 75 \\
 & 123$\to$456 & 19 & 5 & 123 & 69 \\
 & 456$\to$42 & 17 & 7 & 126 & 76 \\
 & 456$\to$123 & 17 & 5 & 126 & 71 \\
\midrule
\multirow{6}{*}{Control @256} & 42$\to$123 & 18 & 6 & 121 & 10 \\
 & 42$\to$456 & 18 & 5 & 121 & 14 \\
 & 123$\to$42 & 19 & 6 & 123 & 13 \\
 & 123$\to$456 & 19 & 5 & 123 & 12 \\
 & 456$\to$42 & 17 & 7 & 126 & 16 \\
 & 456$\to$123 & 17 & 4 & 126 & 14 \\
\bottomrule
\end{tabular}
\end{table}

Cross-seed fresh-attempt correctness for natural-stop errors (26--44\%, pooled 36.1\%) is comparable to the control (24--41\%, pooled 32.4\%), indicating most of this improvement comes from sampling variation rather than budget increase. In contrast, ceiling-hit correctness rises from 8--13\% (control, pooled 10.8\%) to 54--61\% (cross-budget, pooled 57.6\%), a +47pp gain attributable to budget extension.

\subsection*{B. Virtual Budget Extraction}

For the reasoning model, we generate once at @8192 and truncate at the token level to obtain results at lower budgets. Because we use greedy decoding (temperature=0, top\_k=1), the truncated outputs are identical to what would be produced with the lower budget directly---this is a consequence of autoregressive generation with deterministic decoding. We verified this equivalence on a 10-sample pilot, confirming exact match in all cases. This method cannot be used for the sampling experiments (\S4), which require independent stochastic generation at each budget.

\subsection*{C. Statistical Analysis Details}

We use independent sample counting for the Ratchet Effect: each sample is counted once at the most informative budget transition (the transition with the most natural-stop errors for that model). Fisher's exact test for comparing recovery rates between natural-stop and ceiling-hit errors. Clopper-Pearson 95\% CIs for recovery rates. Per-sample bootstrap 95\% CIs for accuracy (10,000 resamples). Cohen's $d$ for effect sizes.

\subsection*{D. Proof of Separability (Per-Model Breakdown)}

The two failure types occupy disjoint regions of the token-count feature space. On 3B at 256 tokens: ceiling-wrong AUC $= 0.815$; natural-wrong AUC $= 0.168$. On 0.5B at 256 tokens: ceiling-wrong AUC $= 0.794$. The natural-wrong vs.\ correct AUC of $0.168$ (below 0.5) arises because natural-wrong answers stop at short lengths while correct answers often run longer, completing their reasoning.

\end{document}


